% ════════════════════════════════════════════════════════════════════════════
% TEMPLATE MAESTRO — GUÍA DE TRABAJO (CLASES PARTICULARES — MATEMÁTICAS)
% ────────────────────────────────────────────────────────────────────────────
% Estética heredada de Guia_Soluciones_Quimicas.tex (Martina, Marzo 2026):
% paleta azul/verde/naranja/rojo, tcolorbox para cajas conceptuales y niveles
% de dificultad en verde/naranja/rojo.
%
% Cómo usar:
%   1. cp guia-base.tex  ../matematicas/2026-S1/semana-NN_tema/Guia_Tema.tex
%   2. Ajustar los \newcommand del bloque "VARIABLES" más abajo.
%   3. Rellenar la PORTADA (contenidos) y las secciones.
%   4. Compilar:  latexmk -pdf Guia_Tema.tex   (o  pdflatex  dos veces)
%
% El .pdf queda fuera de git (ver .gitignore); el .tex sí se versiona.
% ════════════════════════════════════════════════════════════════════════════

\documentclass[11pt,a4paper]{article}
\usepackage[utf8]{inputenc}
\usepackage[T1]{fontenc}
\usepackage[spanish]{babel}
\usepackage{geometry}
\usepackage{amsmath,amssymb}
\usepackage{xcolor}
\usepackage{tcolorbox}
\usepackage{enumitem}
\usepackage{fancyhdr}
\usepackage{titlesec}
\usepackage{multirow}
\usepackage{array}
\usepackage{booktabs}
\usepackage{tikz}
\usepackage{hyperref}

\geometry{margin=2cm, top=2.5cm, bottom=2.5cm}

% ─── VARIABLES DE LA GUÍA ───
% Reemplazar estos valores al crear una guía nueva.
\newcommand{\guiaTitulo}{TÍTULO DE LA GUÍA}
\newcommand{\guiaSubtitulo}{Subtítulo breve}
\newcommand{\guiaUnidad}{Unidad X}
\newcommand{\guiaCurso}{2° Medio --- Matemáticas}
\newcommand{\guiaAlumno}{Alumno}
\newcommand{\guiaSemana}{Semana NN}
\newcommand{\guiaSemestre}{2026-S1}

% ─── Colores ───
\definecolor{azulprincipal}{HTML}{2980B9}
\definecolor{azuloscuro}{HTML}{1A5276}
\definecolor{verdequimica}{HTML}{1ABC9C}
\definecolor{verdeclaro}{HTML}{D5F5E3}
\definecolor{naranja}{HTML}{E67E22}
\definecolor{rojo}{HTML}{E74C3C}
\definecolor{amarillorecuerda}{HTML}{FFF3CD}
\definecolor{amarilloborde}{HTML}{FFC107}
\definecolor{grisfondo}{HTML}{F5F5F5}

% ─── Cajas ───
\tcbset{
  recuerda/.style={
    colback=amarillorecuerda, colframe=amarilloborde,
    fonttitle=\bfseries, title={\color{black} RECUERDA},
    boxrule=1pt, arc=2mm, left=5pt, right=5pt
  },
  propiedad/.style={
    colback=verdeclaro!50, colframe=verdequimica!80!black,
    fonttitle=\bfseries, boxrule=0.8pt, arc=2mm, left=5pt, right=5pt
  },
  formula/.style={
    colback=grisfondo, colframe=azulprincipal,
    boxrule=0.8pt, arc=2mm, left=5pt, right=5pt
  },
  definicion/.style={
    colback=blue!5, colframe=azulprincipal,
    fonttitle=\bfseries, boxrule=0.8pt, arc=2mm, left=5pt, right=5pt
  },
  ejemplo/.style={
    colback=orange!5, colframe=naranja,
    fonttitle=\bfseries, boxrule=0.8pt, arc=2mm, left=5pt, right=5pt
  }
}

% ─── Encabezados ───
\pagestyle{fancy}
\fancyhf{}
\fancyhead[C]{\small\color{gray} \guiaTitulo{} --- \guiaCurso{}}
\fancyfoot[C]{\small\color{gray} Página \thepage}
\renewcommand{\headrulewidth}{0.4pt}

% ─── Títulos ───
% Helper para la barra azul de sección (evita #1 anidado dentro de \parbox,
% que rompe en versiones recientes de titlesec).
\newcommand{\tituloseccion}[1]{%
  \noindent\colorbox{azulprincipal}{%
    \makebox[\dimexpr\textwidth-2\fboxsep\relax][l]{%
      \color{white}\normalfont\Large\bfseries~SECCIÓN \thesection: #1%
    }%
  }%
}
\titleformat{\section}
  {}{}{0em}
  {\tituloseccion}
\titleformat{\subsection}
  {\normalfont\large\bfseries\color{azulprincipal}}
  {\thesubsection}{0.5em}{}

% ─── Niveles de dificultad ───
% Usamos \ (espacio forzado) en lugar de ~, porque ~ es carácter activo
% en babel-español dentro de argumentos de comandos como \colorbox.
\newcommand{\nivelbasico}{%
  \par\smallskip\colorbox{green!60!black}{\color{white}\small\bfseries\ NIVEL BÁSICO\ }\par\smallskip}
\newcommand{\nivelintermedio}{%
  \par\smallskip\colorbox{naranja}{\color{white}\small\bfseries\ NIVEL INTERMEDIO\ }\par\smallskip}
\newcommand{\nivelavanzado}{%
  \par\smallskip\colorbox{rojo}{\color{white}\small\bfseries\ NIVEL AVANZADO\ }\par\smallskip}

% ─── Líneas de respuesta ───
\newcommand{\lineasrespuesta}[1]{%
  \par\vspace{2mm}%
  \foreach \i in {1,...,#1}{\noindent\rule{\textwidth}{0.3pt}\par\vspace{5mm}}%
  \vspace{2mm}%
}

\begin{document}

% ════════════════════════
% PORTADA
% ════════════════════════
\thispagestyle{empty}
\vspace*{3cm}
\begin{center}
  {\Huge\bfseries\color{azulprincipal} GUÍA DE TRABAJO}\\[8pt]
  {\LARGE\bfseries\color{azuloscuro} \guiaTitulo}\\[15pt]
  {\color{azulprincipal}\rule{8cm}{2pt}}\\[15pt]
  {\Large \guiaUnidad{} \textbar{} \guiaCurso}\\[8pt]
  {\large Alumno/a: \guiaAlumno}\\[4pt]
  {\large \guiaSemana{} --- \guiaSemestre}\\[25pt]
  {\itshape\color{gray}
  \guiaSubtitulo}
\end{center}

\vspace{2cm}
\noindent\textbf{Contenidos:}
\begin{enumerate}[leftmargin=2cm]
  \item Tema / subtema 1
  \item Tema / subtema 2
  \item Tema / subtema 3
\end{enumerate}
\newpage

% ════════════════════════
% SECCIÓN 1 — PLACEHOLDER
% ════════════════════════
\section{TÍTULO DE LA SECCIÓN}

\begin{tcolorbox}[recuerda]
Bloque de recuerdos previos / conceptos que deben estar frescos antes de partir.
\end{tcolorbox}

\begin{tcolorbox}[definicion, title=Definición clave]
Texto de la definición.
\end{tcolorbox}

\subsection{Subtema}

Desarrollo teórico breve.

\begin{tcolorbox}[formula]
\[
  \text{fórmula aquí}
\]
\end{tcolorbox}

\begin{tcolorbox}[ejemplo, title=Ejemplo resuelto]
Paso a paso.
\end{tcolorbox}

\subsection*{Ejercicios}

\nivelbasico
\begin{enumerate}[leftmargin=1cm]
  \item Ejercicio básico 1.
  \lineasrespuesta{2}
  \item Ejercicio básico 2.
  \lineasrespuesta{2}
\end{enumerate}

\nivelintermedio
\begin{enumerate}[leftmargin=1cm, resume]
  \item Ejercicio intermedio 1.
  \lineasrespuesta{3}
\end{enumerate}

\nivelavanzado
\begin{enumerate}[leftmargin=1cm, resume]
  \item Ejercicio avanzado 1.
  \lineasrespuesta{4}
\end{enumerate}

\newpage

% ════════════════════════
% RESPUESTAS SELECCIONADAS
% ════════════════════════
\section*{Respuestas seleccionadas}
\addcontentsline{toc}{section}{Respuestas seleccionadas}

\begin{tcolorbox}[propiedad, title=Soluciones]
\begin{itemize}[nosep]
  \item Ej. 1:
  \item Ej. 2:
  \item Ej. 3:
\end{itemize}
\end{tcolorbox}

\end{document}
